\documentclass[10pt,a4paper]{amsart}
\usepackage[margin=19mm]{geometry}
\usepackage{amsmath,amssymb,mathtools}
\usepackage[scr=rsfs,cal=euler]{mathalfa}
\usepackage{xfrac}
\usepackage[unicode,hidelinks,bookmarks=false]{hyperref}
\newcommand{\ie}{i.e.\ }
\newcommand{\cf}{cf.\ }
\newcommand{\resp}{respectively\ }
\newcommand{\Subsection}{\subsection}
\providecommand{\nobreakdash}{\nobreak\hbox{-}}
\setlength{\emergencystretch}{2em}
\multlinegap0pt
\usepackage{amssymb}
\newtheorem{theorem}{Theorem}
\def\del{\partial}
\title{Franklin's identity for $n$-color partitions\\ and companion Beck-type identities}
\author{Cristina Ballantine, Roberto Tauraso}
\date{Source: arXiv:2607.27753v1 (CC BY 4.0)}
\begin{document}
\maketitle

\noindent\emph{Source and licence.} Franklin's identity for $n$-color partitions\\ and companion Beck-type identities. Version arXiv:2607.27753v1 (CC BY 4.0), distributed by arXiv under the Creative Commons Attribution 4.0 International licence, http://creativecommons.org/licenses/by/4.0/, as recorded in the arXiv licence metadata for this exact version and shown on the abstract page https://arxiv.org/abs/2607.27753v1. Full licence notice: \texttt{licenses/arxiv-2607.27753-CC-BY-4.0.txt}.\par\smallskip

\noindent\textit{Editorial scope.} Counted unit: the article's theorem B2 (source0.tex lines 295--301). The article's complete original proof of it is delivered with it (source0.tex lines 303--378, one \texttt{proof} environment of 76 lines). No mathematics is written, repaired or re-derived: the statement and the proof are the article's own lines, in the article's own notation, and no retained line is substituted or re-flowed.  No further source-local context is needed: every cross-reference of every retained line resolves inside the two delivered spans.  Nothing was omitted from any retained span, so the package carries no omission record.  No retained line contains a citation, so the wrapper carries no bibliography and no bibliography entry.  No retained line contains a figure environment, an image-inclusion command or drawing code, so no image is delivered and the package declares no asset dependency.  Editorial formatting only, in addition: an AMS \texttt{amsart} preamble that restates the article's own declarations and macros used by the retained lines, namely the environments \texttt{theorem} on separate counters, together with the standard packages those lines need.  The retained environments print in this document as Theorem 1 in order of appearance, whereas the article prints the same items with the numbering of its own full document.  The complete native source of the article as distributed in the arXiv e-print for this exact version is bundled unchanged as \texttt{sources/206367/source0.tex}.
\begin{theorem}\label{B2}
For $r\geq2$ and $m,j\geq0$, we have
\begin{equation}\label{eqnB2G}
{T}_{j+1,r}(m)-{T}_{j,r}(m)=\sum_{\lambda\in\mathcal{D}_{j,r}(m)}\bar{\ell}(\lambda)-\sum_{\lambda\in\mathcal{O}_{j,r}(m)}\bar{\ell}(\lambda),
\end{equation}
 where $\bar{\ell}(\lambda)$ is the number of distinct parts in $\lambda$.
\end{theorem}

\begin{proof}[Analytic Proof] Let $r\geq 2$. For $k\geq 0$, denote by $\widetilde{\mathcal{O}}_{j,r}(k,m)$ and $\widetilde{\mathcal{D}}_{j,r}(k,m)$ the subsets of partitions with $k$ distinct parts in $\mathcal{O}_{j,r}(m)$ and $\mathcal{D}_{j,r}(m)$, respectively.
We introduce the trivariate generating functions  
$$\widetilde{\mathcal{O}}_r(z,w,q):=
\sum_{m=0}^\infty\sum_{k=0}^\infty\sum_{j=0}^\infty|\widetilde{\mathcal{O}}_{j,r}(k,m)|z^kw^jq^m,$$ and 
$$\widetilde{\mathcal{D}}_r(z,w,q):=\sum_{m=0}^\infty\sum_{k=0}^\infty\sum_{j=0}^\infty|\widetilde{\mathcal{D}}_{j,r}(k,m)|z^kw^jq^m.$$ 
A direct expansion leads to
\begin{align*}
\widetilde{\mathcal O}_r(z,w,q)
&=\prod_{m=1}^\infty(1+wz\sum_{i=1}^{\infty}(q^{rm})^i)^m\cdot
\prod_{m=1}^\infty(1+z\sum_{i=1}^{\infty}(q^{rm})^i)^{(r-1)m}\cdot\prod_{\substack{m=1\\ r\nmid m}}^\infty(1+z\sum_{i=1}^{\infty}(q^{m})^i)^{m}   \\  
&=\prod_{m=1}^\infty\left(1+\frac{wzq^{rm}}{1-q^{rm}}\right)^m\cdot 
\prod_{m=1}^{\infty}\left(1+\frac{zq^{rm}}{1-q^{rm}}\right)^{-m}
\cdot \prod_{m=1}^{\infty}\left(1+\frac{zq^{m}}{1-q^{m}}\right)^{m}\\
&=\prod_{m=1}^{\infty}\left(\frac{1-(1-wz)q^{rm}}{1-(1-z)q^{rm}}\right)^m
\cdot \prod_{m=1}^{\infty}\left(\frac{1-(1-z)q^{m}}{1-q^{m}}\right)^{m},
\end{align*} 
and  
\begin{align*}
\widetilde{\mathcal D}_r(z,w,q)
&=\prod_{m=1}^\infty\Big(1+z\sum_{i=1}^{r-1}(q^{m})^i+wz\sum_{i=r}^{\infty}(q^{m})^i\Big)^m\\
&=\prod_{m=1}^\infty\Big(1+\frac{z(q^m-q^{rm})}{1-q^m}+\frac{wzq^{rm}}{1-q^m}\Big)^m\\
&=\prod_{m=1}^{\infty}\left(\frac{1-(1-z)q^m-(1-w)zq^{rm}}{1-q^{m}}\right)^m.
\end{align*} 
In a similar spirit, we let ${\mathcal T}_{j,r}(k,m)$ be the subset of partitions in $\mathcal{D}_{j,r}(m)$ with exactly $k$ different parts with multiplicity greater than $r$ and less than $2r$. Then
\begin{align*}
{\mathcal T}_r(z,w,q)&:=
\sum_{m=0}^\infty\sum_{k=0}^\infty\sum_{j=0}^\infty|{\mathcal T}_{j,r}(k,m)|z^kw^jq^m\\
&=\prod_{m=1}^\infty\Big(\sum_{i=0}^{r-1}(q^{m})^i+wq^{rm}+wz\sum_{i=r+1}^{2r-1}(q^{m})^i+w\sum_{i=2r}^{\infty}(q^{m})^i\Big)^m\\
&=\prod_{m=1}^\infty\Big(\frac{1-q^{rm}}{1-q^m}+wq^{rm}+\frac{wz(q^{(r+1)m}-q^{2rm})}{1-q^m}+\frac{wq^{2rm}}{1-q^m}\Big)^m\\
&=\prod_{m=1}^{\infty}\left(\frac{1-(1-w)q^{rm}-w(1-z)(q^{(r+1)m}-q^{2rm})}{1-q^{m}}\right)^m.
\end{align*} 
Setting $z=1$ in these three generating functions, it is easy to check that  
\begin{equation}\label{ODT1}
\widetilde{\mathcal O}_r(1,w,q)
=\widetilde{\mathcal D}_r(1,w,q)
={\mathcal T}_r(1,w,q)
=\prod_{m=1}^{\infty}\left(\frac{1-(1-w)q^{rm}}{1-q^{m}}\right)^m.
\end{equation}
We have 
$$\sum_{m=0}^\infty\sum_{j=0}^\infty \Big(\sum_{\lambda\in\mathcal{D}_{j,r}(m)}\bar{\ell}(\lambda)-
\sum_{\lambda\in\mathcal{O}_{j,r}(m)}\bar{\ell}(\lambda)\Big)w^jq^m 
= \left.\frac{\del}{\del z}\right|_{z=1}\!\!\!\!\!\!\bigl(\widetilde{\mathcal D}_r(z,w,q)-\widetilde{\mathcal O}_r(z,w,q)\bigr).$$
Using logarithmic differentiation, we obtain
\begin{align*}\left.\frac{\del}{\del z}\right|_{z=1}\!\!\!\!\!\!\widetilde{\mathcal O}_r(z,w,q)
&=\widetilde{\mathcal O}_r(1,w,q)\sum_{m=1}^\infty m\left(
\frac{wq^{rm}}{1-(1-w)q^{rm}}-q^{rm}+q^{m}\right)
\end{align*}
and, similarly,
\begin{align*}
\left.\frac{\del}{\del z}\right|_{z=1}\!\!\!\!\!\!\widetilde{\mathcal D}_r(z,w,q)
&= \widetilde{\mathcal D}_r(1,w,q)
\sum_{m=1}^\infty \frac{m(q^m-(1-w)q^{rm})}{1-(1-w)q^{rm}}.
\end{align*}
Subtracting these two expressions and applying identity \eqref{ODT1}, we get
\begin{align}\label{RHSB2G}
\sum_{m=0}^\infty\sum_{j=0}^\infty \Big(\sum_{\lambda\in\mathcal{D}_{j,r}(m)}\bar{\ell}(\lambda)&-
\sum_{\lambda\in\mathcal{O}_{j,r}(m)}\bar{\ell}(\lambda)\Big)w^jq^m\nonumber\\
&=\widetilde{\mathcal O}_r(1,w,q)
\sum_{m=1}^{\infty} 
\left(\frac{m(q^m-q^{rm})}{1-(1-w)q^{rm}}-m(q^m-q^{rm})\right)\nonumber\\
&=(1-w)\,\widetilde{\mathcal O}_r(1,w,q)
\sum_{m=1}^\infty \frac{m(q^{(r+1)m}-q^{2rm})}{1-(1-w)q^{rm}}.
\end{align}
On the other hand, evaluating the derivative of ${\mathcal T}_r(z,w,q)$ at $z=1$ gives
\begin{align*}
\sum_{m=0}^\infty\sum_{j=0}^\infty T_{j,r}(m)w^jq^m
&=\left.\frac{\del}{\del z}\right|_{z=1}\!\!\!\!\!\!{\mathcal T}_r(z,w,q)={\mathcal T}_r(1,w,q)
\sum_{m=1}^{\infty}\frac{mw(q^{(r+1)m}-q^{2rm})}{1-(1-w)q^{rm}},
\end{align*}
which implies 
\begin{align}\label{LHSB2G}
\sum_{m=0}^\infty\sum_{j=0}^{\infty} ({T}_{j+1,r}(m)-{T}_{j,r}(m)) w^jq^m=(1-w)\,{\mathcal T}_r(1,w,q)
\sum_{m=1}^\infty \frac{m(q^{(r+1)m}-q^{2rm})}{1-(1-w)q^{rm}}.
\end{align}
Consequently, in view of \eqref{ODT1}, \eqref{RHSB2G} coincides with \eqref{LHSB2G}, which concludes the proof.
\end{proof}

\end{document}
